\documentclass[11pt]{article}
\usepackage[margin=2.4cm]{geometry}
\usepackage{fontspec}
\setmainfont{DejaVu Sans}
\usepackage{juliamono}
\usepackage{fvextra}
\usepackage{xcolor}
\usepackage{xurl}
\usepackage{hyperref}
\usepackage{amssymb}
\newcommand{\checkbox}{$\square$}
\newcommand{\answerchoices}{\checkbox\ YES \quad \checkbox\ NO \quad \checkbox\ CANNOT DETERMINE}
\newcommand{\decisionexplanation}{%
  \par\smallskip\textbf{If NO, give the reason for the mismatch (required). If
  CANNOT DETERMINE, state what context or evidence is missing:}\par
  \rule{\linewidth}{0.4pt}\par\vspace{7pt}\rule{\linewidth}{0.4pt}\par}
\setlength{\parindent}{0pt}\setlength{\parskip}{5pt}
\setlength{\emergencystretch}{2em}
\DefineVerbatimEnvironment{Vsrc}{Verbatim}{fontsize=\small,frame=single,
  breaklines=true,breakanywhere=true}
\DefineVerbatimEnvironment{Vlean}{Verbatim}{fontsize=\scriptsize,frame=single,
  numbers=left,numbersep=5pt,breaklines=true,breakanywhere=true}
\hypersetup{colorlinks=true,urlcolor=blue}

\begin{document}
\begin{center}{\LARGE Proof Verification Document}\\[2pt]{\large reddit\_askmath\_k33fyn}\\[2pt]{\small Reddit proof-check posts (newly confirmed)}\end{center}
\begin{center}\fcolorbox{blue!60!black}{blue!4}{\parbox{0.92\linewidth}{\textbf{Make three independent decisions.} Follow the evidence chain: original source $\longrightarrow$ challenged claim $\longrightarrow$ proposed counterexample $\longrightarrow$ Lean theorem. A record is valid only when all three links are correct. Complete the human sign-off before reading the automated review appendix; those earlier results are background, not an answer key.}}\end{center}
\textbf{Use only this document, except for the public source link.} Everything needed for Questions 2 and 3 is included here. For Question 1, open the public source link when one is shown. If there is no public link, the complete stored problem and solution are included here. If the information is still not enough, choose CANNOT DETERMINE.
For each question, choose \textbf{YES}, \textbf{NO}, or \textbf{CANNOT DETERMINE}. Choose CANNOT DETERMINE if the document does not provide enough information for you to decide; do not mark YES if you are uncertain. If you find a specific mismatch, choose NO and explain it in the space below that question. Do not rewrite the claim, add assumptions, or substitute a different counterexample.
\textbf{How copied text is shown.} The problem, quoted passage, and proposed counterexample come from the stored record. The Lean theorem statement and Lean files are copied exactly from the certificate. Some long lines continue on the next printed line; this is only page formatting and does not change the copied text. Line numbers are added only to help you find passages. If the stored record contains a character that cannot be printed, the document shows \texttt{[U+XXXX CONTROL]} in its place.
\begin{description}
\item[Challenged claim.] The exact statement from the source that this record says is false.
\item[Proposed counterexample.] A concrete case that should satisfy the claim's assumptions but make its conclusion false.
\item[Lean certificate.] A Lean file proving a formal theorem. Successful compilation proves the theorem as written, not that it is the right translation of the source.
\end{description}
\section*{Question 1 --- Does this proof verification document accurately represent the source?}
\textbf{What to decide.} Compare the challenged passage below with the original source---either the public link or the complete source text in this document. Answer YES only if they make the same claim and include the same assumptions or conditions. If this document changes the meaning, leaves out a condition, or points to a different passage, answer NO.
\subsection*{Original source}
\textbf{Source record:}
\begin{Vsrc}
qedbench_certificate_availability_triage/data/records_reddit_exhaustive.jsonl
\end{Vsrc}
\textbf{Public source:}\par
\url{https://www.reddit.com/r/askmath/comments/k33fyn/proofs_using_unique_prime_factorizations/}\par
\subsection*{Problem or question from the source (verbatim)}
\begin{Vsrc}
Proofs using unique prime factorizations.
\end{Vsrc}
\subsection*{Challenged passage recorded from the source (stored text)}
\begin{Vsrc}
Suppose x is an odd integer, and suppose its prime factorization is x = p\_1 · p\_2 · p\_3 · ... · p\_n.
\end{Vsrc}
\subsection*{Surrounding source lines (stored text with added line labels)}
The labels refer to lines in the solution text stored with the source record; the source-line content is not transliterated.
\begin{Vsrc}
0001 | I've just started a proofs course. I was told we're allowed to use the Fundamental Theorem of Arithmetic, that each integer greater than 1 can be expressed as the product of primes in exactly one way.
0002 | 
0003 | I understand this fact, but I'm not sure how exactly to use it in proofs. This is what I came up with, but it seems more complicated than necessary. I'm wondering if there's a simpler way.  
0004 | 
0005 |   
0006 | 
0007 | Prove: If x is an odd integer, then x^(4) is an odd integer.
0008 | 
0009 | Proof: Suppose x is an odd integer, and suppose its prime factorization is x = p\_1 · p\_2 · p\_3 · ... · p\_n. Since x is odd, p\_1, p\_2, p\_3, ... , p\_n are all not equal to 2. x^(4) = p\_1^(4) · p\_2^(4) · p\_3^(4) · ... · p\_n^(4). Since the prime factorization of x^(4) does not contain 2, x^(4) is odd. QED
0010 | 
0011 | &#x200B;
0012 | 
0013 | I know you could prove the same result by defining an odd number to be of the form 2n+1 witn n ∈ ℕ, but I just used this as an example because I thought it was something I could probably prove with the Fundamental Theorem. I just want to make sure I am using it correctly.
0014 | 
0015 | Thank you :)
\end{Vsrc}
\textbf{A YES means:} the proof verification document faithfully represents both the source context and the challenged claim, without dropping a condition or changing its meaning.
\section*{Question 2 --- Does the proposed counterexample disprove that claim?}
\textbf{Claim to test (repeated for convenience).}
\begin{Vsrc}
Suppose x is an odd integer, and suppose its prime factorization is x = p\_1 · p\_2 · p\_3 · ... · p\_n.
\end{Vsrc}
\subsection*{How the record says the argument uses this claim}
This is the record's description, not an established fact; check it against the source.
\begin{Vsrc}
The proof may apply the stated prime-factorization theorem to every odd integer x and obtain the displayed nonempty product of primes.
\end{Vsrc}
\subsection*{Proposed counterexample (exact values or construction)}
\begin{Vsrc}
Restrict the theorem’s application to positive odd integers x>1 (equivalently x>=3); handle units separately.
\end{Vsrc}
\subsection*{Why the record says the counterexample works}
Check this explanation independently; it is evidence to assess, not an answer key.
\begin{Vsrc}
The hypotheses allow every odd integer, including 1, but 1 has no nonempty prime factorization. Thus the stated theorem cannot be applied as claimed without excluding units (and, depending on conventions, handling negative integers separately).
\end{Vsrc}
\subsection*{How to decide}
\begin{enumerate}
\item List every assumption that applies to the challenged claim.
\item Check that the proposed example satisfies every one of those assumptions.
\item Recompute the claimed outcome and the actual outcome using exact reasoning.
\item Confirm that the conflict is with the challenged claim itself, not an unrelated sentence.
\end{enumerate}
\textbf{A YES means:} the example is fully specified, satisfies all relevant assumptions, and makes the challenged claim false. If a necessary value, condition, or calculation is missing, choose CANNOT DETERMINE.
\section*{Question 3 --- Does the final Lean theorem statement verify this same counterexample?}
\textbf{What to decide.} A Lean file can compile while proving the wrong theorem. Compare these three links independently:
\begin{enumerate}
\item The Lean definitions represent the same mathematical objects, assumptions, and operations as the challenged claim.
\item The Lean construction uses the same concrete values or example shown in Question 2.
\item The final theorem statement asserts that the challenged claim fails on that example, not merely a weaker or nearby statement.
\end{enumerate}
\subsection*{Selected final Lean theorem statement (verbatim UTF-8 excerpt)}
This is an exact contiguous slice of the certificate file, from the theorem keyword through the original \texttt{:=} token where its proof begins. It is shown first so the formal target is not buried in the complete file.
\begin{Vlean}
theorem reddit_askmath_k33fyn_refutation :
    OddInt 1 ∧
    ¬ Factorized 1 ∧
    ¬ FTAApplicable 1 ∧
    ¬ (∀ x : Int, OddInt x → Factorized x) :=
\end{Vlean}
\subsection*{Project-local imported Lean files (complete, verbatim UTF-8)}
This section follows the certificate's import commands recursively and prints every repository-owned Lean source file on which the final theorem statement depends. Each file is copied in full, including its definitions and controls. External Mathlib, Std, and Lean Init source is not duplicated.
No project-local imported Lean files were found. External Mathlib, Std, and Lean Init source is not duplicated in this proof verification document.
\subsection*{Complete counterexample Lean file (verbatim UTF-8)}
The complete Lean file is included below so you can check every definition and supporting result used by the final theorem statement. Use it together with the project-local imported Lean files above.
\textbf{Lean certificate file:}
\begin{Vsrc}
partB_feasibility/reddit_certificates/reddit_askmath_k33fyn.lean
\end{Vsrc}
\begin{Vlean}
/-!
Record: reddit_askmath_k33fyn
Public source: https://www.reddit.com/r/askmath/comments/k33fyn/proofs_using_unique_prime_factorizations/

Exact target scope: the source proof applies a prime-factorization theorem stated
for integers greater than `1` to every odd integer `x`, and therefore assumes a
displayed nonempty product of primes for every such `x`.  The concrete witness
is the allowed odd integer `x = 1`: its factor list is empty, with empty product
`1`, while every nonempty list of positive natural primes has product at least
`2`.  The independently verified negative control is `x = 3`, represented by
the nonempty factor list `[3]`.

Modelling judgments: `x` ranges over `Int`; oddness means `x = 2*k + 1` for an
integer `k`; factors range over `Nat`; natural divisibility is witnessed by a
natural quotient; and a natural prime has lower bound `2` and no divisors other
than `1` and itself.  List multiplication is recursive with empty product `1`,
and `Factorized` explicitly requires a nonempty all-prime list whose natural
product is embedded into `Int` with `Int.ofNat`.  Thus this model is the
positive-prime product displayed by the source, not a signed factorization and
not a finite prime lookup table.  `FTAApplicable x` records only the theorem's
material side condition `1 < x`: it excludes the unit witness and accepts the
odd control `3`; no full fundamental-theorem-of-arithmetic result is asserted.
The units and negative odd integers require separate handling rather than being
silently included in this positive factorization model.
-/

namespace RedditAskmathK33fyn

def OddInt (x : Int) : Prop :=
  ∃ k : Int, x = 2 * k + 1

def NatDivides (d n : Nat) : Prop :=
  ∃ q : Nat, n = d * q

def PrimeNat (p : Nat) : Prop :=
  2 ≤ p ∧ ∀ d : Nat, NatDivides d p → d = 1 ∨ d = p

def product : List Nat → Nat
  | [] => 1
  | p :: ps => p * product ps

def AllPrime : List Nat → Prop
  | [] => True
  | p :: ps => PrimeNat p ∧ AllPrime ps

def Factorized (x : Int) : Prop :=
  ∃ ps : List Nat,
    ps ≠ [] ∧ AllPrime ps ∧ x = Int.ofNat (product ps)

def FTAApplicable (x : Int) : Prop :=
  1 < x

theorem product_pos_of_allPrime :
    ∀ ps : List Nat, AllPrime ps → 0 < product ps
  | [], _ => by decide
  | p :: ps, hall => by
      exact Nat.mul_pos
        (Nat.lt_of_lt_of_le (by decide : 0 < 2) hall.1.1)
        (product_pos_of_allPrime ps hall.2)

theorem nonempty_allPrime_product_ge_two
    (ps : List Nat) (hne : ps ≠ []) (hall : AllPrime ps) :
    2 ≤ product ps := by
  cases ps with
  | nil => exact False.elim (hne rfl)
  | cons p rest =>
      have hp : 2 ≤ p := hall.1.1
      have hrest : 0 < product rest :=
        product_pos_of_allPrime rest hall.2
      have hp_le_product : p ≤ p * product rest := by
        have hmul := Nat.mul_le_mul_left p hrest
        simpa using hmul
      exact Nat.le_trans hp hp_le_product

theorem positive_le_three_cases
    (n : Nat) (hpos : 0 < n) (hle : n ≤ 3) :
    n = 1 ∨ n = 2 ∨ n = 3 := by
  cases n with
  | zero => simp at hpos
  | succ n =>
      cases n with
      | zero => exact Or.inl rfl
      | succ n =>
          cases n with
          | zero => exact Or.inr (Or.inl rfl)
          | succ n =>
              cases n with
              | zero => exact Or.inr (Or.inr rfl)
              | succ n => simp at hle

theorem primeNat_three : PrimeNat 3 := by
  constructor
  · decide
  · intro d hd
    rcases hd with ⟨q, hproduct⟩
    have hdpos : 0 < d := by
      cases d with
      | zero => simp at hproduct
      | succ d => exact Nat.zero_lt_succ d
    have hqpos : 0 < q := by
      cases q with
      | zero => simp at hproduct
      | succ q => exact Nat.zero_lt_succ q
    have hdle : d ≤ 3 := by
      calc
        d = d * 1 := by simp
        _ ≤ d * q := Nat.mul_le_mul_left d hqpos
        _ = 3 := hproduct.symm
    have hqle : q ≤ 3 := by
      calc
        q = q * 1 := by simp
        _ ≤ q * d := Nat.mul_le_mul_left q hdpos
        _ = d * q := Nat.mul_comm q d
        _ = 3 := hproduct.symm
    rcases positive_le_three_cases d hdpos hdle with hd | hd | hd
    · subst d
      exact Or.inl rfl
    · subst d
      rcases positive_le_three_cases q hqpos hqle with hq | hq | hq
      · subst q
        simp at hproduct
      · subst q
        simp at hproduct
      · subst q
        simp at hproduct
    · subst d
      exact Or.inr rfl

theorem oddInt_one : OddInt 1 := by
  refine ⟨0, ?_⟩
  decide

theorem not_factorized_one : ¬ Factorized 1 := by
  intro hfactor
  rcases hfactor with ⟨ps, hne, hall, heq⟩
  have hge : 2 ≤ product ps :=
    nonempty_allPrime_product_ge_two ps hne hall
  have hproduct : product ps = 1 := by
    have habs := congrArg Int.natAbs heq
    simpa using habs.symm
  rw [hproduct] at hge
  exact (by decide : ¬ (2 ≤ 1)) hge

theorem refutes_published_universal_step :
    ¬ (∀ x : Int, OddInt x → Factorized x) := by
  intro published
  exact not_factorized_one (published 1 oddInt_one)

/-- Load-bearing refutation of the literal universal proof step: `1` satisfies
the source's odd-integer premise, fails the required nonempty positive-prime
factorization, and is outside the stated theorem's `1 < x` applicability scope.
-/
theorem reddit_askmath_k33fyn_refutation :
    OddInt 1 ∧
    ¬ Factorized 1 ∧
    ¬ FTAApplicable 1 ∧
    ¬ (∀ x : Int, OddInt x → Factorized x) := by
  refine ⟨oddInt_one, not_factorized_one, ?_,
    refutes_published_universal_step⟩
  change ¬ ((1 : Int) < 1)
  decide

/-- Meaningful negative control: the same oddness and positive-factorization
encoding accepts `3`, with the genuine nonempty prime list `[3]`, and `3` is
inside the corrected applicability scope. -/
theorem reddit_askmath_k33fyn_negative_control :
    OddInt 3 ∧ PrimeNat 3 ∧ Factorized 3 ∧ FTAApplicable 3 := by
  constructor
  · refine ⟨1, ?_⟩
    decide
  constructor
  · exact primeNat_three
  constructor
  · refine ⟨[3], ?_, ?_, ?_⟩
    · decide
    · exact ⟨primeNat_three, True.intro⟩
    · decide
  · change (1 : Int) < 3
    decide

end RedditAskmathK33fyn

\end{Vlean}
\subsection*{Proof-safety information}
The automated scan found no unfinished proofs, new axiom declarations, unsafe declarations, external code hooks, or uses of Lean's compiled decision procedure. This scan concerns proof construction only; it does not establish that the theorem matches the source claim.
\textit{Technical scan counts:} native\_decide=0, sorry=0, axiom\_decl=0, unsafe=0, extern=0. \textit{Technical tier:} KERNEL-ONLY.
\textbf{A YES means:} Lean represents the same objects and assumptions, uses the same proposed counterexample, and the final theorem statement asserts the same contradiction. Successful compilation by itself is not enough.
\clearpage
\section*{Human sign-off}
Complete this block from your own review of Questions 1--3. Do not use the automated appendix below as an answer key.\par
\textbf{Reviewer name:} \rule{9cm}{0.4pt}\par
\textbf{Date:} \rule{4cm}{0.4pt}\par
\medskip\textbf{1. Does the proof verification document accurately represent the source and challenged claim?}\par
\answerchoices\par
\decisionexplanation
\medskip\textbf{2. Does the proposed counterexample satisfy the assumptions and make that claim false?}\par
\answerchoices\par
\decisionexplanation
\medskip\textbf{3. Does the final Lean theorem statement represent and verify that same counterexample?}\par
\answerchoices\par
\decisionexplanation
\medskip\textbf{Overall decision:}\par
\checkbox\ VALIDATED (all three answers are YES)\par
\checkbox\ REJECTED (at least one answer is NO)\par
\checkbox\ NEEDS ANOTHER REVIEWER (at least one CANNOT DETERMINE and no NO)\par
\medskip\textbf{Additional comments.} Use this space for any extra source, counterexample, calculation, or Lean details that support your decisions.\par\rule{\linewidth}{0.4pt}\par\vspace{7pt}\rule{\linewidth}{0.4pt}
\clearpage
\section*{Appendix --- Prior automated review (background only)}
\textbf{These results are not an answer key.} They come from the earlier Opus/Fable verification campaign and are printed only so an auditor can see what was previously checked. The human decisions above must be made independently. Raw status codes are retained in small type for traceability.
\textbf{Overall automated disposition:} Cleared by automated screening; an independent search also found a counterexample.
\par\textit{Audit code:} \nolinkurl{PASS_CORROBORATED}
\subsection*{Independent search with no access to the stored example}
\textbf{Result:} An independent search found a counterexample without seeing the stored example.
\par\textit{Audit code:} \nolinkurl{WITNESS_FOUND}
\par\textbf{Cases checked:}
\begin{Vsrc}
All 502 odd integers x with -501 <= x <= 501. For each x, complete recursive enumeration of every non-decreasing list of primes of length n >= 1 whose product is exactly x (and, for the sign audit, exactly |x|). The enumeration is exhaustive because any prime factor of x divides |x| and hence lies in [2,|x|], and any product of n primes has absolute value >= 2^n, so n <= log2(|x|). (Shortened here; the complete note remains in the machine-check JSON.)
\end{Vsrc}
\subsection*{Independent re-check of the proposed counterexample}
\textbf{Result:} A separate checker recomputed the proposed counterexample and confirmed it.
\par\textit{Audit code:} \nolinkurl{CONFIRMED}
\par\textbf{Checker's note:}
\begin{Vsrc}
Verified from scratch. Nothing under validation/machine_checks/searches/ and no Lean file was read. The only file consulted was the input record. The primary witness x = 1 survives every convention shift I tested; the claim as quoted genuinely fails on it.
\end{Vsrc}
\subsection*{Independent review of the source-to-Lean correspondence}
\textbf{Result:} A separate reviewer judged that Lean proves the same failure using the same example.
\par\textit{Audit code:} \nolinkurl{MATCHES}
\par\textbf{Theorem summary from that reviewer:}
\begin{Vsrc}
The load-bearing theorem RedditAskmathK33fyn.reddit_askmath_k33fyn_refutation literally asserts a four-part conjunction about the integer 1: (i) OddInt 1, i.e. there is an integer k with 1 = 2*k+1 (k = 0), so 1 satisfies the source's 'x is an odd integer' premise; (ii) NOT Factorized 1, where Factorized x means there is a NONEMPTY list ps of naturals, every entry PrimeNat (2 <= p and every natural divisor of p is 1 or p), with x = Int.ofNat (product ps) -- i.e. (Shortened here; the complete note remains in the machine-check JSON.)
\end{Vsrc}
\par\textbf{Difference noted by that reviewer:}
\begin{Vsrc}
None load-bearing; all three links hold. (1) Definitions: OddInt (exists k, x = 2k+1) over Int correctly captures 'x is an odd integer' and does include the unit 1; PrimeNat is the standard nat primality (2 <= p, only divisors 1 and p, with NatDivides d n := exists q, n = d*q); product is the recursive list product with empty product 1; Factorized requires a NONEMPTY all-prime list, which is the faithful reading of the quoted display 'x = p_1 * p_2 *... (Shortened here; the complete note remains in the machine-check JSON.)
\end{Vsrc}
\subsection*{Follow-up review}
No follow-up review was needed or recorded.
\textit{This appendix is only background. You do not need any separate machine-output files to answer the three questions.}
\end{document}